\section{Propuesta}

\subsection{Módulo de Geometría Solar}
Para poder determinar de manera correcta la carga térmica que cada una de las paredes del espacio habitacional que se estudia tienen, es necesario considerar la presencia del sol. Como ya hemos visto, es posible calcular con precisión la posición del sol en un determinado día y hora del año. A través de cálculos geométricos es posible saber qué paredes son las que se encuentran más expuestas a los efectos del astro. \\

Usando la técnica de trazado de rayos, se pueden identificar las caras que se encuentran expuestas directamente a los haces de radiación solar. A continuación se muestra una figura que ejemplifica este concepto.

\begin{figure}[htbp]
	\centering
	\includegraphics[width=0.7\textwidth, angle=0]{images/sunraysboxmodel.pdf}
	\caption{Modelo de caja expuesta a los haces de radiación solar.}
\end{figure}